\documentclass[final,a4,xcolor={svgnames,dvipsnames}]{beamer}
\usetheme{Madrid}
\title{le projet RefPerSys}
\author{Basile STARYNKEVITCH}
\date{mai 2025}
\usepackage{relsize}
\usepackage{luacode}
\usepackage[svgnames]{xcolor}
\usepackage{alltt}
%\usepackage{wasysym}
\usepackage{hyperref}
%\usepackage{newunicodechar}

 \hypersetup{
   colorlinks   = true, %Colours links instead of ugly boxes
   urlcolor     = DarkBlue, %Colour for external hyperlinks
   linkcolor    = DarkGreen, %Colour of internal links
   citecolor   = DarkMagenta, %Colour of citations
   frenchlinks = true,
 }

\begin{document}
 \begin{luacode*}
   local gitpip=io.popen("git log --no-color --format=oneline -1 --abbrev=16 --abbrev-commit -q | cut -d' ' -f1")
   gitid=gitpip:read()
   gitpip:close()
 \end{luacode*}
 \newcommand{\mygitid}{\luadirect{tex.print(gitid)}}

 \begin{frame}
   \titlepage

   \small{(un moteur d'inférences, ``intelligence artificielle'' symbolique)}

   \bigskip

   \begin{center}
   \begin{relsize}{-2}
     8 rue de la Faïencerie\\
     92340 Bourg-la-Reine
   \end{relsize}

   \bigskip

   \textbf{les opinions me sont personnelles}
   
   \bigskip
   
     \href{mailto:basile@starynkevitch.net}{\texttt{basile{@}starynkevitch.net}}

     \bigskip

     \begin{fbox}
       {Journées Du Logiciel Libre - Lyon - mai 2025}
     \end{fbox}
     

     git {\texttt \mygitid}
   \end{center}
   
 \end{frame}

 \begin{frame}{Plan}

   \begin{relsize}{-1.5}
   \tableofcontents
   \end{relsize}
   
   \bigskip

   \begin{center}
     \relsize{+2}{\href{http://refpersys.org/}{\texttt{refpersys.org}}}

     \bigskip
     \includegraphics[width=0.45\textwidth]{refpersys-logo}

     \smallskip

     \relsize{-2}{Logo par Gaëtan Tapon}
   \end{center}
 \end{frame}


 %%%%%%%%%%%%%%%%%%%%%%%%%%%%%%%%%%%%%%%%%%%%%%%%%%%%%%%%%%%%%%%%
 \section{Introduction}
 \label{sec:intro}
 
 \begin{frame}{motivations pour RefPerSys}

   \begin{itemize}
   \item souveraineté numérique \small{(moins dépendre de logiciels
     états-uniens, même libres)}
   \item s'occuper 
   \item expérimenter un logiciel
   \item populariser et affiner des idées méconnues (celles de
     \href{https://fr.wikipedia.org/wiki/Jacques_Pitrat}{\it Jacques
       Pitrat})
   \item \textbf{un projet \textcolor{red}{en cours de dévelopement} qui cherche des contributeurs}
   \end{itemize}

   \bigskip
   \textbf{Le code source est libre} sous licence GPLv3+/CeCILLv2
   \textbf{sans garantie} sur
   \href{https://github.com/RefPerSys/RefPerSys}{\texttt{github.com/RefPerSys/RefPerSys}}
   et peut-être même difficilement utilisable en mai 2025.

   \bigskip

   \begin{center}
     \large\textsc{\textcolor{red}{REF}lexive \textcolor{red}{PER}sistent \textcolor{red}{SYS}tem}
   \end{center}
 \end{frame}
 
 \begin{frame}{historique et systèmes inspirants pour RefPerSys}

   \begin{itemize}

   \item \textbf{les moteurs d'inférences des systèmes experts} des années 1990
     \href{https://clipsrules.net/}{\textsc{CLIPSrules}} (NASA)

   \item \textbf{le système \textcolor{red}{CAIA}} (``chercheur
     artificiel en intelligence artificielle'') de J.Pitrat en
     {\texttt{\href{https://github.com/bstarynk/caia-pitrat}{github.com/bstarynk/caia-pitrat/}}}
     qui auto-génère son demi-million de lignes de C (1990-2018)
     
   \item \textbf{les outils d'aide à la démonstration} ou assistants
     de preuve \href{https://rocq-prover.org/}{\textsc{Rocq}} (ex
     \textsc{Coq})

     \item La problématique de \textbf{\textcolor{red}{l'amorçage des
         compilateurs}}: tous les compilateurs se compilent
       eux-mêmes. Idéalement \textit{\textcolor{red}{RefPerSys devrait générer son
         propre code}} mais en génère une petite partie en mai 2025.
   \end{itemize}

   \bigskip

   En France, il existe une grosse culture mathématique de systèmes
   formels de preuve. Ma philosophie est plus pragmatique et plus ludique.
   
 \end{frame}
 
 \begin{frame}{une propriété méconnue de Linux: les greffons à foison}


   Un programme sous Linux est {\relsize{-1}{(très souvent)}}
   utilisateur de nombreuses bibliothèques dynamiques
   {\relsize{-1}{(``shared objects'')}}, les fichiers
   \texttt{lib*.so}. Ce mécanisme optimise l'utilisation de la mémoire
   RAM et il est utilisé pour les
   \href{https://fr.wikipedia.org/wiki/Plugin}{greffons}
        {\relsize{-1}{(``plugins'')}} chargés par la fonction
        \textbf{\texttt{\href{https://man7.org/linux/man-pages/man3/dlopen.3.html}{dlopen(3)}}}. Le
        programme obtient par
        \textbf{\texttt{\href{https://man7.org/linux/man-pages/man3/dlsym.3.html}{dlsym(3)}}},
        durant l'exécution, les fonctions d'un greffon à partir de
        leur nom.

   \medskip
     
   Un grand nombre de programmes sous Linux {\relsize{-1.5}{(dont le
       navigateur \href{https://firefox.org}{\textsc{Firefox}},
       l'éditeur programmable
       \href{https://www.gnu.org/software/emacs/}{\textsc{GNU emacs}},
       le compilateur \href{https://gcc.gnu.org/}{\textsc{Gcc}}, les
       interprètes de commandes
       \href{https://www.gnu.org/software/bash/}{\textsc{Bash}} ou
       \href{https://zsh.org}{\textsc{Zsh}} ou
       \href{https://lua.org}{\textsc{Lua})}}} acceptent des greffons,
   pour permettre d'en étendre le comportement. Ceux-ci peuvent être
   codés par d'autres dévelopeurs.

   \bigskip

   En pratique \textbf{un programme qui s'exécute} (son processus)
   \textbf{peut charger \textcolor{red}{un grand nombre} de greffons}
   (dizaines ou centaines de milliers).

   Donc \textbf{un programme \textcolor{red}{peut générer beaucoup de greffons}
     temporaires}.

   \smallskip

   \small{Cf
   \href{https://github.com/bstarynk/misc-basile/blob/master/manydl.c}{\texttt{manydl.c}
     dans \texttt{github.com/bstarynk/misc-basile}}.}

 \end{frame}

 \begin{frame}{Chargement et déchargement de code, une idée ancienne}

   Dans les années 1950 ou 1960, \textbf{les ordinateurs avaient peu
     de mémoire} {\relsize{-1.5}{(quelques kilo-octets ou un-demi
       méga-octet)}} \textbf{et}
   {\relsize{-0.5}{l'\href{https://fr.wikipedia.org/wiki/IBM_1130_et_1800}{IBM
         1130}, les
       \href{https://fr.wikipedia.org/wiki/CAB500}{CAB500} et CAB2020
       de la
       \href{https://fr.wikipedia.org/wiki/Soci\%C3\%A9t\%C3\%A9_d\%27\%C3\%A9lectronique_et_d\%27automatisme}{SEA},
       l'\href{https://fr.wikipedia.org/wiki/IBM_360_et_370}{IBM 360}
       et d'autres machines}}\footnote{Les µ-ordinateurs des années
   1970 ou 1980 aussi:
   \href{https://fr.wikipedia.org/wiki/IBM_PC/XT}{PC/XT}, TRS-80,
   \href{https://fr.wikipedia.org/wiki/Amiga_500}{Amiga 500}.}
   \textbf{chargeaient le code par bouts}. \textbf{Un segment de code
     devenu inutile était supprimé de la mémoire centrale}, et la
   gestion fine du chargement et déchargement de ces segments
   nécessitait la coopération des programmeurs humains, du système
   d'exploitation (le ``moniteur'') et était une source de bogues
   intéressants.

   \bigskip

   Sous Linux, il n'est guère indispensable de décharger un
   greffon\footnote{L'exception serait un code hypothétique de calcul
   qui tourne pendant des mois sur un \textit{super-calculateur} du
   \href{https://top500.org/}{\texttt{top500.org}} en chargeant des
   millions de greffons générés temporairement.\\} par
   \textbf{\texttt{\href{https://man7.org/linux/man-pages/man3/dlclose.3.html}{dlclose(3)}}}. La
   puissance des machines et la performance du noyau
   {\relsize{-1.5}{(sous-système de gestion de la mémoire virtuelle)}}
   permet de ne pas le faire\footnote{On accepte alors une fuite
   mémoire sans dommage.\\}.

 \end{frame}

 \begin{frame}{Génération de code, une idée ancienne}

  Les compilateurs traduisent un code source (compréhensible par le
  dévelopeur) en du code machine\footnote{La notion de code objet ou
  intermédiaire est apparue progressivement dans les années 1960;
  certaines architectures -dont Harvard, qu'on retrouve encore dans
  les micro-contrôleurs à l'intérieur de nos claviers et de nos
  chaudières- séparent la mémoire du code et celles des données.} et
  existaient déjà dans les années 1950 ou 1960 (Cobol, Fortran,
  \href{https://fr.wikipedia.org/wiki/Programmation_automatique_des_formules}{PAF},
  \href{https://en.wikipedia.org/wiki/Information_Processing_Language}{IPL-V},
  \href{https://fr.wikipedia.org/wiki/Lisp}{Lisp}).

  \bigskip

  De nos jours, \textbf{les dévelopeurs utilisent des logiciels}
  spécialisés \textbf{générateurs de code} (en C par exemple, ou en
  Java ou Ocaml): \href{https://www.gnu.org/software/bison/}{GNU
    bison} comme \href{https://carburetta.com/}{Carburetta} ou
  \href{https://gallium.inria.fr/~fpottier/menhir/}{Menhir} génèrent
  tous du code pour faciliter l'analyse lexicale ou syntaxique. La
  théorie des languages formels est mûre
  (\href{https://fr.wikipedia.org/wiki/Noam_Chomsky}{Chomsky}).
 \end{frame}

 \begin{frame}{Génération de code, une idée moderne}

  Des bibliothèques logicielles pour générer du code existent:
  \href{https://www.gnu.org/software/lightning/}{GNU lightning},
  \href{https://asmjit.com/}{AsmJit}, et le compilateur
  \href{https://gcc.gnu.org/}{GCC} fournit sa
  \href{https://gcc.gnu.org/onlinedocs/jit/}{libgccjit}

  \bigskip
  
  Le Common Lisp libre \href{https://sbcl.org/}{SBCL} génère du code
  machine à chaque interaction de son utilisateur. C'est un langage \href{https://fr.wikipedia.org/wiki/Homoiconicité}{homoïconique}: le code y est manipulable comme donnée.

  \bigskip
  
  La plupart des implémentations de Java (cf
  \href{https://openjdk.org/}{OpenJDK}) ont une
  \href{https://fr.wikipedia.org/wiki/Machine_virtuelle_Java}{JVM} qui
  traduit petit à petit le code-octet (java bytecode) en du code
  machine.

 \end{frame}

 \begin{frame}{Évaluation partielle: une manière d'optimiser}

Beaucoup de calculs longs\footnote{La météo nationale calcule le temps
de demain en des heures de calcul sur super-calculateur; de même en
astrophysique -simulation de collisions de galaxies-; ou en
bioinformatique -simulation des mouvements moléculaires dans une
cellule-.\\} (durant des heures ou des semaines) dépendent de plusieurs
paramètres.
\medskip

Certains paramètres sont constants durant tout le calcul\footnote{Pour
la météo, les mesures physiques entrées des paramètres des capteurs,
mais les résultats internes intermediaires varient durant ce calcul
itératif.\\}.
\medskip

Dans certains cas, il peut être utile de particulariser le code pour ces paramètres constants.

Par exemple simpliste, une affectation pourrait être \texttt{u := a*x
  + b*y + c} et si on sait que \texttt{a} et \texttt{c} sont nuls elle
se simplifie en \texttt{u := b*y} plus rapide à calculer.


 \end{frame}

 \begin{frame}{Évaluation partielle: formalisation}
   
Formellement:

\begin{itemize}
\item on connait le code source $\sigma$ (millions de lignes)
\item ce code calcule sur des paramètres fixes \textcolor{DarkGreen}{$\pi_1, \pi_2, \ldots \pi_n$}
\item ce code calcule sur des entrées variables $v_1, v_2, \ldots v_p$
  
\item les résultats obtenus sont $r_1, r_2 \ldots r_q = \sigma (\textcolor{DarkGreen}{\pi_1, \pi_2, \ldots \pi_n}, v_1, v_2 \ldots v_p)$
\item par des algorithmes complexes on produit et génère un code
  dérivé $\delta$ qui dépend des paramètres \textcolor{DarkGreen}{$\pi_i$} tel que
  $\delta(v_1, v_2, \ldots v_p) = \sigma(\color{DarkGreen}{\pi_1, \pi_2, \ldots \pi_n}\color{black}{,  v_1, v_2 \ldots v_p)}$
\end{itemize}

La dérivation du code $\delta$ à partir du code source $\sigma$ et des
valeurs connues des paramètres fixes \textcolor{DarkGreen}{$\pi_i$}
est compliquée (peut être indécidable).
 \end{frame}


 %%%%%%%%%%%%%%%%%%%%%%%%%%%%%%%%%%%%%%%%%%%%%%%%%%%%%%%%%%%%%%%%
 \section{Formaliser, représenter, traiter les connaissances}
 \label{sec:form-conn}
 
 \begin{frame}{Variété des connaissances}
   Un programme est écrit par un dévelopeur\footnote{Sans présupposé
   sur son genre, étant constaté qu'il y a moins de femmes que
   d'hommes qui pratiquent ou étudient la programmation, et que le
   métier d'informaticien me semble hélàs plus masculin que dans les années
   1980!\\} qui a besoin de connaître:
   \begin{relsize}{-0.5}
   \begin{itemize}
   \item des connaissances sur l'algorithmie (les tris, les automates, les arbres, \ldots)
   \item des connaissances sur les langages de programmation et les systèmes d'exploitation utilisés
   \item des connaissances sur le problème qu'il doit résoudre dans le logiciel qu'il développe.
   \item des connaissances sur les ordinateurs\footnote{Un
   super-calculateur est différent d'un portable, même si tous deux
   utilisent Linux!\medskip\\} sur lequel s'exécutera le logiciel.
   \item des connaissances sur les contraintes (délais, séverité des
     bogues) et habitudes du projet: le logiciel de pilotage d'un
     métro automatique a d'autres contraintes qu'un logiciel
     bureautique ou un compilateur.
   \end{itemize}
   \end{relsize}
 \end{frame}


 \begin{frame}{Formaliser -ou pas- les connaissances.}

   Certaines connaissances ne sont pas ou ne peuvent pas être
   formalisées. Les legislations\footnote{Cf
   \href{https://www.cnil.fr/fr/reglement-europeen-protection-donnees}{\it RGPD}
   ou
   \href{https://artificialintelligenceact.eu/high-level-summary/}{\it EU-AI
     act}}, les habitudes, les intérêts individuels, la science et la
   technique informatiques imposent des limitations.

   \bigskip
   
   Le dévelopeur explicite certaines connaissances dans le code source et
   dans sa documentation:
   \begin{relsize}{-0.5}
   \begin{itemize}
   \item choix du (ou des) langages (formels) utilisés et de la forme des entrées et sorties.
   \item importance des noms choisis
   \item importance des commentaires
   \item séparer l'interface logicielle de son implantation
   \item importance de l'historique du code source et du projet
     logiciel et des ``normes'' ou ``standards''
   \item cohérence relative et imparfaite entre le code source et sa
     documentation\footnote{La \href{https://fr.wikipedia.org/wiki/Programmation_lettrée}{programmation lettrée} est difficile,
     et souvent coûteuse, peu utile ou inpraticable.\medskip\\}
   \end{itemize}
   \end{relsize}
 \end{frame}

 \begin{frame}{méta-connaissances}
   Ce sont les connaissances sur les connaissances.


   \begin{itemize}
   \item pour formaliser des connaissances: par exemple grammaire
     formelle de C ou de la syntaxe BNF des grammaires formelles.
   \item pour utiliser des connaissances: par exemple: un compilateur
     C\footnote{Les compilateurs C n'arrivent guère à compiler en
     l'optimisant une fonction C de cent milles lignes de code dans un
     fichier C généré par un logiciel.\medskip\\}peut supposer que
     chaque fichier C a un nombre raisonnable de définition de
     fonctions de taille raisonnable.
   \end{itemize}

   \bigskip

   Un dévelopeur aguerri a des \href{https://theses.fr/1990PA066799}{méta-connaissances sur la programmation}.
   \bigskip
   
   \textbf{Le même formalisme peut être utilisé pour représenter des
     méta-connaissances qui peuvent s'appliquer à elles-mêmes.}
     
 \end{frame}

 \begin{frame}{connaissances déclaratives ou procédurales}

   \textbf{Une \textcolor{red}{connaissance déclarative} est un énoncé qui peut être utilisé de
   plusieurs façons}. Par exemple: (grammaire française)
   \textit{``l'adjectif s'accorde en genre et en nombre avec le nom
     auquel il se rapporte''}. Elle s'étend facilement (grammaire
   latine: \textit{``\ldots\! et en cas''}). Elle peut être formalisée:
   $\forall x \in \mathbb{R}^{*}, \hspace{1.5em} 0x = 0  \hspace{1.em} \wedge \hspace{1.em} x^0 = 1 \hspace{1.em}  \wedge \hspace{1.em} 0+x = x$

   \bigskip

   \textbf{Une \textcolor{red}{connaissance procédurale} n'est
     utilisable que d'une seule façon}.
   Le code binaire\footnote{La
   complexité et le coût des microprocesseurs s'explique aussi par le
   fait que le silicium réordonne dynamiquement l'exécution des
   instructions machines (``out of order superscalar processors'').}
   est  procédural. En C
{\relsize{-0.5}{\texttt{\textbf{while} (x>0) printf(\textquotedbl\%d\textbackslash{n}\textquotedbl,x--);}}}


   \smallskip

En fait \textbf{le distinguo est graduel} et ces définitions
varient selon les époques et les lieux
\footnote{Aux USA
 ``declarative programming'' a un sens différent de ``programmation
 déclarative'' en France!\\}.
%%
%%   \bigskip
 \end{frame}

 \begin{frame}{Avantages de connaissances plus déclaratives}
     
   \textbf{Fournir des connaissances formalisées et plus déclaratives
     augmente la productivité des dévelopeurs}.

   \bigskip

   Le génie logiciel (``software engineering'') tend à privilégier
   cette productivité\footnote{Au détriment d'autres facteurs, y
   compris énergétiques: par exemple l'utilisation de PHP ou Python dans des
   serveurs Web consomme plus de courant et de temps que le même
   service codé en C et compilé. Le numérique consommerait 5 à 10\% de la
   production électrique européenne.\medskip}.

   \bigskip
   
   \textbf{Les outils logiciels pour les dévelopeurs doivent favoriser
     la déclarativité des formalismes.}

   \bigskip

   L'ordre dans lequel sont formalisées les connaissances déclaratives n'a pas d'importance!
 \end{frame}



 %%%%%%%%%%%%%%%%%%%%%%%%%%%%%%%%%%%%%%%%%%%%%%%%%%%%%%%%%%%%%%%%
 \begin{frame}{Traitement en machine des connaissances plus déclaratives}

   Le formalisme déclaratif est représenté pour son utilisateur par du
   texte (ou des images) formel\footnote{La présentation de ce texte -
   couleurs et polices de caractères - en facilite la lisibilité et sa
   compréhension par des humains.}, et en machine par des structures
   de données compliquées. C'est mathématiquement un graphe ou un
   hypergraphe. On suppose qu'il tient en mémoire centrale.
   
   Un formalisme plus ou moins déclaratif peut être ``interprété'' ou
   ``compilé'' par l'outil logiciel l'acceptant.

   \begin{itemize}
   \item \textbf{interprétation}: l'outil logiciel parcourt
     directement et itérativement, pour traiter ce formalisme,
     l'hypergraphe de sa représentation et modifie les données
     traitées.
     
   \item \textbf{compilation}: l'outil logiciel transforme cette
     representation interne en d'autres représentations plus
     procédurales, in fine interprétées ou transformées en du code
     exécutable (ou plus efficace mais moins déclaratif).
     
   \end{itemize}
   
 \end{frame}

 \begin{frame}{une définition de l' ``Intelligence Artificielle''}

   Enseigné par Jacques Pitrat:
   
   \begin{block}{définition de l'IA}
     \textbf{Un logiciel d'``intelligence artificielle'' est un programme qui peut surprendre \textcolor{red}{agréablement} ses utilisateurs}.
   \end{block}

   \bigskip
   
   \textbf{Tout gros logiciel est bogué} {\relsize{-1}{(et son nombre de bogues reste non-nul)}} \textbf{donc surprend \textcolor{red}{désagréablement} ses utilisateurs}.

   \bigskip

   Un logiciel d'IA peut être partiellement introspectif ou
   \href{https://fr.wikipedia.org/wiki/Réflexion_(informatique)}{réflexif}:
   en particulier en inspectant son code, ses données et ses
   \href{https://fr.wikipedia.org/wiki/Pile_d'exécution}{piles
     d'exécution}.

   \bigskip
   
\begin{quote}{\href{https://fr.wikipedia.org/wiki/Loi_de_Hofstadter}{Loi de Hofstadter}} Il faut toujours plus de temps que prévu, même en tenant compte de la loi de Hofstadter.
\end{quote}

Elle s'applique au dévelopement de RefPerSys et à l'écriture de ces
transparents.
   
 \end{frame}

 \begin{frame}{Règles d'inférences  et bases de règles}

   Une \textbf{règle d'inférences} comporte:

   \begin{itemize}
   \item un ensemble (\textit{non ordonné}, fini) de conditions
     (déclanchant la règle)
   \item une séquence d'actions (pouvant s'exécuter quand ou si les
     conditions sont déclanchées)
   \item des méta-informations supplémentaires
   \end{itemize}

   Exemple: Si $x$ est le père de $y$ et si $z$ est la fille de $y$
   alors on peut en déduire que $x$ est le grand-père de $z$.

   \bigskip
   

   Une \textbf{base de règles} contient un ensemble (\textit{non
     ordonné}, fini) de règles d'inférences, et des méta-informations
   supplémentaires
 Une \textbf{base de faits} contient un ensemble
 (\textit{non ordonné}, fini) de faits (expressions supposées vraies, entités, \ldots)

 \bigskip
 
   Un \textbf{système expert} comporte une ou des bases de règles et
   des bases de faits
 \end{frame}
 
 \begin{frame}{Persistance}

   Un logiciel
   \href{https://fr.wikipedia.org/wiki/Persistance_(informatique)}{persistant}
   sauvegarde son état sur disque et peut redémarrer (le lendemain?)
   avec cet état.

   \bigskip

   \textbf{RefPerSys persiste son état dans des fichiers textuels} (JSON et
   code C++ généré) \textbf{portables}.

   \bigskip

   Comment gérer la co-évolution du logiciel et des données
   persistentes et du code?
   
 \end{frame}

%
%
%%%%%%%%%%%%%%%%%%%%%%%%%%%%%%%%%%%%%%%%%%%%%%%%%%%%%%%%%%%%%%%%
 
 \section{RefPerSys en mai 2025}
 \label{sec:refpersys-mai2025}
 \begin{frame}{RefPerSys en mai 2025}
   \begin{itemize}
   \item dévelopé sous Linux (Debian ou Ubuntu ou Mint, parfois Raspbian)
   \item c'est du logiciel libre sous licence GPL/CeCILL
   \item historique accessible par \href{https://github.com/RefPerSys/RefPerSys/}{\texttt{github.com/RefPerSys/RefPerSys/}}
   \item codé en C++17 et compilé par un \href{https://gcc.gnu.org/}{GCC} récent (v14 ou v15)
   \item nombreuses dépendances (de logiciels libres) dont
       \href{https://fltk.org}{FLTK} et
       \href{https://github.com/ianlancetaylor/libbacktrace}{libbacktrace}
       et \href{https://carburetta.com/}{Carburetta}
     \item quelques contributeurs bénévoles (Suède, Inde, Egypte,
       France, USA?)
   \end{itemize}

   \bigskip
   
   Pour le compiler: définir la variable d'environment
   \texttt{REFPERSYS\_TOPDIR} au répertoire source (obtenu par
   \texttt{git clone}) puis \texttt{make config} (logiciel interactif)
   puis \texttt{make -j4}

   Votre \texttt{\$HOME/tmp/} doit exister (votre répertoire temporaire)
 \end{frame}
 %
 \begin{frame}{tester RefPerSys}
   
   Quelques tests:

   \begin{itemize}
     \item \texttt{make test-load} et d'autres (essayez
      \texttt{make showtests}).

      \item \texttt{make redump} ou \texttt{make
     altredump} pour la persistance et regénération (doit être suivi
      de \texttt{make clean} puis \texttt{make -j4})

    \item \texttt{make testcarb1} (à améliorer)
   \end{itemize}
 \end{frame}
 \begin{frame}{Données manipulées par RefPerSys}

   \begin{itemize}
   \item le pointeur \texttt{nullptr} a un status spécial (une absence de valeur).
   \item les valeurs {\relsize{-0.5}{(un mot de 64 bits souvent pointeur, en C++ \texttt{Rps\_Value})}} peuvent être:
     \begin{enumerate}
     \item entiers directs (62 bits)
     \item des valeurs scalaires immuables (chaînes, flottants, \ldots)
     \item des valeurs composites immuables (ensembles, tuples,
     fermetures,  instances)
       \item des reférences (en C++ \texttt{Rps\_ObjectRef}) à des objets lourds et mutables
     \end{enumerate}
     \item les valeurs composites et les objets peuvent être
       temporaires ou persistents
       \item chaque valeur a sa classe (modèle
         \href{https://en.wikipedia.org/wiki/ObjVlisp}{ObjVLisp})
         et son type  {\relsize{-0.5}{(en C++ \texttt{enum class Rps\_Type})}}
       \item quelques objets sont des racines préexistantes (fichier
         C++ généré \href{https://github.com/RefPerSys/RefPerSys/blob/master/generated/rps-roots.hh}{\texttt{generated/rps-roots.hh}})
   \end{itemize}

 \end{frame}

 \begin{frame}{Cadres d'appel explicites}
   
   Le code C++ représente un cadre d'appel réifié par
   \texttt{Rps\_CallFrame}. On déclare en C++ ses valeurs locales par une macro \texttt{RPS\_LOCALFRAME}, e.g.

   \begin{relsize}{-1}
   \begin{alltt}
     RPS\_LOCALFRAME(RPS\_ROOT\_OB(\_1aGtWm38Vw701jDhZn)\footnote{Cet
     objet décrit le cadre d'appel pour en faciliter la réflexion et
     l'inspection.\medskip\\}, //the\_agenda,\\
     \hspace{8em}            RPS\_NULL\_CALL\_FRAME, // no caller frame\\
     \hspace{8em}            Rps\_ObjectRef obtasklet;\\
     \hspace{8em}            Rps\_InstanceValue descrval;\\
     \hspace{8em}            Rps\_ClosureValue clostodo;\\
                );
   \end{alltt}
   \end{relsize}
%   \smallskip
   
   En C++ le cadre d'appel courant est \texttt{\_} et (dans l'exemple ci-dessus) la référence locale à son objet \texttt{obtasklet} se note \texttt{\_f.obtasklet}, etc.
   
   Ces cadres d'appels RefPerSys sont parcourus par le ramasse-miettes et les routines d'introspection.

   On peut en C++ ajouter des collections C++ de valeurs au cadre
   d'appel courant en explicitant le code C++ qui les parcourt.
 \end{frame}

 
 \begin{frame}{Les valeurs scalaires dans RefPerSys}
   En plus des entiers directs (``tagged integers'') elles peuvent être une zone mémoire copiable par le ramasse-miettes pour (au choix):

   \begin{itemize}
     \item les chaînes de caractères emboîtées UTF-8 
     \item les flottants emboîtés
     \item les valeurs JSON
   \end{itemize}

   \bigskip
  
  \textbf{On devrait pouvoir rajouter d'autres types de valeurs}
  (images, vecteurs de nombres complexes, \ldots) \textbf{et en
    \textcolor{red}{générer le code} C++ les gérant}.
 \end{frame}
 
 \begin{frame}{Les objets dans RefPerSys}
   \textbf{Ils ont chacun un identifiant ``mondialement''
     \textcolor{red}{unique} et aléatoire} et persistant, leur
   \textbf{objid}. \relsize{-0.5}{En mémoire 128 bits, dans les fichiers textuels
   (JSON) ou le C++ une séquence telle que
   \texttt{\_1Io89yIORqn02SXx4p} à partir duquel on forme des noms de fonctions C ou C++ tels que \texttt{rpsapply\_61pgHb5KRq600RLnKD}}.

   \medskip
   
   \begin{relsize}{-0.5}
   De plus chaque objet contient en propre:
   \begin{itemize}
   \item son verrou C++ \texttt{std::recursive\_mtx}
     \item son espace de persistance (une référence à un objet ou \texttt{nullptr} pour les objets temporaires)
     \item sa classe, qui est un \texttt{Rps\_ObjectRef} RefPerSys persistant.
     \item une association d'attributs (en C++ un
       \texttt{std::map}). Chaque attribut est un objet, sa valeur
       quelconque (sauf \texttt{nullptr})
     \item une séquence de composants (en C++ un \texttt{std::vector}) de valeurs (parfois \texttt{nullptr})
       \item une charge utile optionnelle (sous classe de
         \texttt{Rps\_Payload}) qui connait son objet propriétaire.
       \item quelques pointeurs de fonctions C++ qui sont
         \href{https://man7.org/linux/man-pages/man3/dlsym.3.html}{\texttt{dlsym}}-ables
         et
         \href{https://man7.org/linux/man-pages/man3/dladdr.3.html}{\texttt{dladdr}}-ables
   \end{itemize}
   \end{relsize}
   \end{frame}

 \begin{frame}{Les valeurs composites dans RefPerSys}

   Elles sont \textbf{\textcolor{red}{immuables}} et
   \textit{copiables} par le ramasse-miettes. Chacune peut être
   temporaire ou persistante.

   \begin{relsize}{-0.5}
   \begin{itemize}
     \item \textbf{ensemble} fini d'objets ordonné par leur \textit{objid}. Le
       test d'appartenance est rapide (par dichotomie).
     \item \textbf{tuple} (ou n-uplet) de réferences d'objets; une composante
         de tuple peut-être absente (\texttt{nullptr}).
       \item \textbf{fermeture}: sa routine est donnée par un objet
           qui décrit cette routine. Ses valeurs closes sont
           immuables. De plus elles peuvent contenir un méta-descripteur (référence d'objet) et un méta-rang (nombre descriptif).
           \item \textbf{instance}: organisée comme une fermeture,
             elles contient une connective (objet), des fils
             (valeurs fixes), et un méta-descripteur et un méta-rang.
   \end{itemize}


   Attention à notre terminologie: \textbf{une instance
     \textcolor{red}{n'est pas} un objet} (car immuable) mais toutes
   les valeurs composites ont leur classe et leur type. Les instances
   pourraient être utiles pour représenter des n{\oe}uds d'une syntaxe
   ou des faits, etc...
   \end{relsize}
 \end{frame}

 \begin{frame}{application d'une fermeture}
   
   Chaque fermeture connait son code (par un objet, sa ``connective'')
   C++ (une fonction C++ dont le nom commence par \texttt{rpsapply}).

   \bigskip
   
   Appliquer une fermeture à des arguments (tous des valeurs), c'est
   appeler cette fonction C++ (avec la valeur fermeture).

   \bigskip
   
   En C++ cette application se fait par du code tel que \texttt{\_f.clostodo.apply1(\&\_, \_f.obtasklet)}

   \bigskip

   Cette application renvoie deux valeurs (dans un couple en C++
   \texttt{Rps\_TwoValues}) qui par convention ne doivent pas être
   toutes deux \texttt{nullptr}. Elle devrait être presque aussi
   rapide qu'un appel de méthode \texttt{virtual} en C++.
   
 \end{frame}
 
 \begin{frame}{envoi de message}
   
   Inspiré par Smalltalk ou Common Lisp, il est plus souple (mais plus
   coûteux en temps) qu'un appel de méthode \texttt{virtual} de C++.

   \smallskip
   On envoie à une valeur cible un message donné par un sélecteur et des arguments.
   En C++ par exemple \texttt{Rps\_TwoValues two = \_f.refpersysv.send3(\&\_, \_f.gencodselob, \_f.dumpdirnamev, \_f.tempsuffixv, \_f.genstoreob);}

   \bigskip
   \textbf{Le selecteur est un objet}

   \bigskip
   L'envoi de message est implémenté par application d'une
   fermeture. Chaque classe RefPerSys est un objet dont la charge
   utile contient une superclasse et un dictionnaire de méthodes
   asssociant à un selecteur la fermeture réalisant l'envoi. Ce
   dictionnaire peut évoluer avec le temps.

   \smallskip
   \begin{small}
     Ce modèle ressemble à celui de JavaScript, mais s'appuie sur les
     classes (pas de prototypes).
   \end{small}
 \end{frame}
 
 \begin{frame}{Le ramasse-miettes dans RefPerSys}
   
   \href{https://fr.wikipedia.org/wiki/Ramasse-miettes_(informatique)}{Il}
   (``garbage collector'' ou glaneur de cellules) est conçu pour
   devenir plus efficace (en C++ \texttt{class Rps\_GarbageCollector})
   donc parfois copieur générationnel:

   \begin{itemize}
   \item les routines du R-M devraient être facilement générables.
   \item les valeurs scalaires ou composites (toutes immuables)
     doivent être copiables par le R-M
   \item les objets lourds sont fixes en mémoire (contenant le verrou \texttt{std::recursive\_mutex} de chaque objet) et le R-M les marquent.
     \item parcours de la pile d'exécution et de ses cadres d'appels (\texttt{Rps\_CallFrame})
   \end{itemize}
 \end{frame}
 
 \begin{frame}{La persistance dans RefPerSys}

   Le code C++ de persistance ressemble par sa forme au R-M. Les deux devraient devenir générés.
   \medskip
   Au démarrage:
   \begin{itemize}
   \item à froid préallocation statiques des zones mémoire des objets
     racines
   \item chargement des objets vides de chaque espace
     \item remplissage des objets et initialisation des pointeurs de fonction conventionnels par 
         \href{https://man7.org/linux/man-pages/man3/dlsym.3.html}{\texttt{dlsym}}
   \end{itemize}

   \bigskip
   \textbf{Le tas persistant} (``persistent heap'') \textbf{est donc \textcolor{red}{chargé} à chaque
   démarrage à partir de fichiers textuels.}
   
   \bigskip
   \textbf{À la fin normale du processus} (et peut-être tous les
   quarts d'heure) \textbf{on écrit sur disque le tas persistant mis à
     jour et le code généré en C++}. Le code C++ utilise
   \href{https://man7.org/linux/man-pages/man3/dladdr.3.html}{\texttt{dladdr}}
   pour retrouver le nom des pointeurs de fonctions.
   
 \end{frame}

 \begin{frame}{Écriture du tas persistant}
   \textbf{Le tas persistant est écrit sur disque}, avec le code C++ généré qui
   va avec.

   \bigskip
   
   Chaque espace de persitance correspond à un fichier JSON. On peut
   donc distinguer le système (les objets nécessaires à la génération
   de code C++) et plusieurs ``applications''.

   \bigskip
   
   \textbf{On pourrait} à l'avenir \textbf{utiliser d'autres
     représentations textuelles} (HJSON?, Yaml?) pour les données
   persistées. \textbf{\textcolor{red}{Il faut veiller à leur
       portabilité}}: un tas écrit sur x86-64 devrait être utilisable
   et chargeable sur \href{https://riscv.org}{RISCV}.
   
 \end{frame}


 \begin{frame}{agenda et millitâches}

   \textbf{\textcolor{red}{L'agenda est un objet racine persistant}}
   \relsize{-0.5}{\texttt{\_1aGtWm38Vw701jDhZn} nommé \texttt{the\_agenda}, unique de
   sa classe \texttt{\_3s7ztCCoJsj04puTdQ}
   \textit{\texttt{agenda$\in$class}}}.

   \bigskip

   \textbf{Cet agenda contient des millitâches} (``tasklets''), qui
   sont des objets (persistants ou temporaires). \textbf{L'exécution
     de chaque millitâche doit durer \textcolor{red}{quelques
       millisecondes}}.

     \textbf{Il existe plusieurs\footnote{Leur nombre est figé au
       démarrage de RefPerSys (entre 3 et 24, selon la puissance du
       processeur et son nombre de c{\oe}urs).} filaments d'exécution}
     (\href{https://fr.wikipedia.org/wiki/Threads_POSIX}{POSIX
       threads}).

 \end{frame}

 \begin{frame}{exécution parallèle des millitâches de l'agenda}

 \textbf{Chacun des filaments d'exécution} de l'agenda
  \textbf{boucle} sur:

     \begin{enumerate}
     \item choix\footnote{En mai 2025 l'agenda contient quelques files
     d'attente de priorité haute, normale, ou basse de millitâches en
     attente d'exécution.} et retrait d'une millitâche à exécuter dans
       l'agenda
     \item exécution de cette millitâche par envoi d'un message
       \item cette exécution est rapide (millisecondes) et peut
         inspecter et modifier l'agenda (peut-être s'y ajouter à nouveau)
     \end{enumerate}

     \bigskip
     
     Cette boucle {\relsize{-0.5}{(interrompue par le ramasse-miettes)}} s'arrête quand
     RefPerSys termine normalement.

     \bigskip
     \textbf{Des événements extérieurs} \footnote{interaction
     graphique ou Web avec l'utilisateur, ou
     \href{https://www.jsonrpc.org/}{JSONRPC}, terminaison d'un
     processus -compilation d'un greffon généré en C++?-, interruption
     temporelle toutes les secondes?}  peuvent créer et ajouter (ou
     ôter) des millitâches dans l'agenda.
   
 \end{frame}
%%%%%%%%%%%%%%%%%%%%%%%%%%%%%%%%%%%%%%%%%%%%%%%%%%%%%%%%%%%%%%%%
 
 \section{Ce qui reste à faire dans RefPerSys}
 \label{sec:refpersys-a-faire}
 \begin{frame}{A faire avec votre aide}
   \begin{itemize}
   \item définir une syntaxe (règles, métarègles, ...) avec des exemples (cf \href{https://github.com/RefPerSys/RefPerSys/blob/master/carbrepl_rps.cbrt}{carbrepl\_rps.cbrt})
   \item accroître le tas persistent pour décrire déclarativement dans
     des objets RefPerSys le code C++ à générer (en reprenant des idées de \href{https://arxiv.org/abs/1109.0779}{GCC MELT})
     \item coder une interface graphique {\relsize{-0.5}{(esquissée avec \href{https://fltk.org}{FLTK} ou autre et son code C++ doit en partie être généré à partir d'objets RefPerSys)}}.
   \item coder les générateurs de codes C++ (y compris l'inférence) et machines\footnote{Le
   code machine généré n'est pas persisté car non-portable; il
   pourrait être généré rapidement par du code C++ généré qui utilise \textit{GNU lightning} ou autre!\medskip\\}
     \item générer du code C++ qui remplace le code C++ écrit à la
       main (qu'on \texttt{git mv} dans \texttt{attic/}
   \end{itemize}

 \end{frame}

 \begin{frame}{co-dévelopement itératif de RefPerSys}
   
   \textbf{Ce développement de RefPerSys devrait être itératif,
     collaboratif, et communautaire}. Le tas persistant de RefPerSys
   doit générer de plus en plus de code (qui remplacerait celui
   existant codé\footnote{J'ai pris soin de nommer les identifiants
   C++ de manière à espérer les générer} à la main) à partir de
   formalismes plus déclaratifs donc plus concis et plus
   compréhensible.

   \bigskip
   
   \textbf{Le formalisme de RefPerSys est appelé à évoluer avec le
     temps en des representations de plus en plus déclaratives} qui
   requièrent votre collaboration.

   \bigskip
   
   \textit{Un système d' ``intelligence artificielle'' doit expliquer
     son comportement aux utilisateurs pour être acceptable.}
 \end{frame}

 \begin{frame}{applications possibles de RefPerSys}

   \begin{itemize}
   \item jeux
   \item optimisation et pilotage de codes de calculs sous Linux
   \item aide à la décision (jardinage, médical?, pilotage de projets?
     CAO?)
   \item et toute autre chose
   \item \ldots
   \end{itemize}

   En respectant le cadre légal et l'éthique du libre.
 \end{frame}
 
\end{document}

%%%%%%%%%%%%%%%%%%%%%%%%%%%%%%%%%%%%%%%%%%%%%%%%%%%%%%%%%%%%%%%%
%% Local Variables: ;;
%% compile-command: "./build-refpersys-jdll2025.sh" ;;
%% End: ;;
%%%%%%%%%%%%%%%%%%%%%%%%%%%%%%%%%%%%%%%%%%%%%%%%%%%%%%%%%%%%%%%%
